\writeSectionFrame{\presentationTitle}{Übung}

\begin{frame}
\frametitle{Ziel}
\begin{itemize}
    \item Implementieren Sie das Singleton-Muster in einer Java-Anwendung.
    \item Verstehen Sie die Anwendung und Einschränkungen des Singleton-Musters.
\end{itemize}
\end{frame}

\begin{frame}
\frametitle{Aufgabenbeschreibung}
\begin{enumerate}
    \item \textbf{Erstellen Sie eine Singleton-Klasse:}
    \begin{itemize}
        \item Implementieren Sie eine Klasse \texttt{Logger}, die nur eine einzige Instanz während der Laufzeit haben soll.
        \item Die Klasse \texttt{Logger} soll eine Methode \texttt{log(String message)} enthalten, die eine Nachricht in der Konsole ausgibt.
    \end{itemize}
    
    \item \textbf{Verwenden Sie den Singleton:}
    \begin{itemize}
        \item Implementieren Sie eine \texttt{main}-Methode oder eine Testklasse, die die \texttt{Logger}-Instanz verwendet, um mehrere Nachrichten zu protokollieren.
        \item Stellen Sie sicher, dass immer nur eine Instanz des \texttt{Logger} verwendet wird, auch wenn mehrere Threads gleichzeitig auf den \texttt{Logger} zugreifen.
    \end{itemize}
\end{enumerate}
\end{frame}